\clearpage
\fancyhead[L]{\small ASTRONOMICAL YEARBOOK.}
\hypertarget{native-610-lower}{}
\begin{center}\small Göttingische gelehrte Anzeigen. Nos. 81 and 82. Pages 801..810. 21 May 1818.\\[3pt]\rule{.72\linewidth}{.3pt}\end{center}

\emph{Astronomical Yearbook for the Year {\upshape 1820}, together with a collection of the latest papers, observations, and reports pertaining to the astronomical sciences. Calculated and published with the approval of the Royal Academy of Sciences by} Dr. I. E. \textsc{Bode}, Royal Astronomer, etc. Berlin. Available from the author and on commission through F. Dümmler. 256 pages in octavo, with one copperplate.

\begingroup\brokenpenalty=10000\relax
The year 1820 is distinguished by several noteworthy astronomical phenomena. Besides the great solar eclipse of 7 September, which will also be annular in our region, there will be several occultations of planets by the Moon. One of Mars on 28 January and two of Jupiter on 4 June and 18 October will be visible in Berlin. --- The supplements open with the astronomical observations made at the Berlin Observatory in 1816. Among these we single out the observation of the solar eclipse of 19 November and that of the lunar eclipse of 4 December. --- Observations in Halberstadt by Captain von \textsc{Wahl}. The latitude of Halberstadt (Moritzplan) was found from solar observations with an 18-inch \textsc{Troughton} reflecting circle to be $51^\circ\,54'\,5''9$. From a solar eclipse and an occultation of Aldebaran, Mr. \textsc{von Wahl} calculates the longitude as $1'\,17''6$ of time east of the Seeberg Observatory. The conjunction times determined here for several places included in the comparison, using Mr. \textsc{von Wahl}'s calculation, differ by a few seconds from the results obtained by others in their calculations. --- Visible variations in the brightness of Algol in the years 1817---1819, calculated by Prof. \textsc{Wurm}. --- Astronomical notices and observations of the opposition of Jupiter in 1816 by the astronomer \textsc{Derflinger} at Kremsmünster. --- Astronomical remarks and calculation of the opposition of Vesta in 1815 and of the solar eclipse of 1820, by Dr. \textsc{Gerling} in Cassel. This article also contains a correction of a hasty judgment about atmospheric refraction in the yearbook for 1819. --- On \textsc{Kepler}'s problem, by State Councillor and Knight \textsc{von Schubert}. In a sense, this paper is a detailed expansion of \textsc{Laplace}'s solution of the problem by means of an infinite series in the sines of multiples of the mean anomaly. Its coefficients in turn take the form of infinite series in successive powers of the eccentricity. In a later supplement the author carried the admittedly very laborious expansion to the 13th power of the eccentricity. Insofar as the expansion is of mathematical interest, one must be grateful to the author for it, even if one judges the astronomical value of a solution by such an infinite series, compared with so-called indirect solutions, quite differently from him and finds that precisely what he seeks to \SourceLeaf{611}advance in favor of the former at the expense of the latter can be applied only in the reverse sense. For it is precisely when using the series that one cannot carry the accuracy beyond the point to which one has worked out the values of the numerical coefficients, whereas the accuracy of the so-called indirect method is determined solely by the accuracy of the sine tables used. In any case, it is well known that the latter, when suitably arranged, are incomparably more convenient for fairly large eccentricities. --- Astronomical observations at the Vienna Observatory in 1816 by \textsc{Triesnecker} and \textsc{Bürg}. Observations of eclipses of Jupiter's satellites, occultations of stars, several planetary oppositions, the lunar eclipse of 9 June, and the solar eclipse of 18 November. --- On adjusting the meridian telescope, by Prof. \textsc{Littrow}. This paper addresses the problem of determining the position of the meridian telescope from observed transits when the errors are of finite magnitude. Strictly speaking, this problem is of more mathematical than astronomical interest, since in practice the only errors that occur are those to be regarded as infinitesimal. Moreover, applying this method depends on knowing the true times of culmination, which can be obtained only with the aid of another instrument. For an independent adjustment of the meridian telescope, one could take only the differences between the true times of culmination as known and must then assume either that the instrument's axis is correctly adjusted or that the inclination is given. Of course, this procedure too is suitable only for meridian telescopes of lesser rank, and an adjustment that depends on knowing the right ascensions is inadmissible for an instrument of the first rank. --- From a letter by the same astronomer, an excerpt is given from an investigation aimed at determining the planetary masses that perturb the Earth's motion. Since the investigation is printed at length in the Zeitschrift für Astronomie, we take the liberty of making a few remarks about it. Comparison with the journal shows that a misprint in the present excerpt made the uncertainty (more accurately, the probable error) of the mass of Venus ten times as large as it should be. The author compared 189 Greenwich solar observations with \textsc{von Zach}'s solar tables and, by combining the observations in groups, derived 43 equations of condition that contain 11 unknown quantities He confines himself, however, to determining by the method of least squares the values of only one of them, namely the corrections to the epoch, the mass of Venus, the mass of Mars, and the mass of the Moon. We do not quite see, however, why the epoch alone among the elliptic elements is considered in preference to the longitude of perigee and the eccentricity. The assertion that it is difficult to specify the epoch with some precision would probably also need correcting. Furthermore, the judgment (Zeitschrift, p. 285) that nutation can be determined from solar observations with greater certainty than the mass of Venus is untenable, at least when the solar longitudes in question are determined from observed right ascensions. For what is observed directly consists of differences in right ascension between the Sun and the fundamental fixed stars, and these differences are only weakly affected by nutation. The solar longitudes themselves already already involve the nutation computed for the fixed stars, and one must therefore always end up merely recovering almost exactly the same nutation that was taken as the basis of the calculation. Finally, as for the results themselves that Mr. \textsc{Littrow} himself obtained for the four quantities mentioned, it would be highly desirable to repeat the calculation, because it contains several visible errors, and it is difficult to decide whether they are misprints, slips of the pen, or errors in calculation. The nature of the method of least squares requires the coefficients of $\mu'''$ in the first equation and of $\mu$ in the third equation on p. 292 to be equal, whereas here they are entirely different. The coefficients of $\mu'''$ in the second equation and of $\mu'$ in the third equation should likewise be exactly the same. As the sum of the squares of the errors before applying \SourceLeaf{612}the correction obtained, Mr. \textsc{Littrow} found 973, and after applying it, 727. It seems to us, however, that this reduction should have been considerably greater. We cannot accept the view that it is more appropriate to add the signed errors themselves than their squares. For the vanishing of that sum attests \emph{only} to the effectiveness of the correction to the epoch and says nothing at all about the other corrections. We hope that the author will take these remarks as a sign of the attention with which we have followed his valuable investigation and, in keeping with his promise, will soon carry it out on a larger scale. --- Astronomical observations in 1816 at the Prague and Vilna observatories. The former include the complete observation of the solar eclipse of 18 November and the opposition of Ceres, and the latter include the opposition of Pallas. --- Professor \textsc{Bessel}'s observations at the Königsberg Observatory are especially valuable. They cover six planetary oppositions, five occultations of stars, the lunar eclipse of 4 December 1816, the zenith distances of the Sun at both solstices, and the right ascension of the Pole Star. --- Remarks on the occasion of the great solar eclipse of 19 November 1816. Since a great solar eclipse is one of the rare occasions when the weather conditions at a large number of places become publicly known, the editor had the happy idea of compiling all the places where a clear sky allowed the eclipse to be observed in full, then those where changeable weather allowed only intermittent glimpses of it, and finally those where a completely overcast sky prevented any part of it from being seen. The result shows that all these places lie scattered in a jumble, without connection or regular pattern, and is therefore well suited to weakening belief in specific weather prophecies to some degree. The author adds a few remarks on the observation that during total eclipses the darkness is usually less pronounced than one expects. What he puts forward as an attempt to explain the ring sometimes seen around the Moon during total eclipses, however, has not become quite clear to us. --- Observations of eclipses of Jupiter's satellites and of occultations of stars at the Greenwich Observatory in 1811---1814, and observations of the great comet of 1811 made there as well. --- New elements of Vesta's orbit by Professor \textsc{Gerling}. --- Astronomical observations at the Göttingen Observatory by Prof. \textsc{Gauss}. The lunar eclipse of 4 December 1816, the altitude of the celestial pole at the new observatory determined from 158 observations of the North Star, and observations of the summer solstice of 1817. These observations yielded the same result that has puzzled astronomers for several years. The altitude of the pole derived from circumpolar stars comes out several arcseconds greater than that derived from observations of the Sun, even though the observations mentioned were made without the counterweight normally fitted to the objective and suspected by some of being a possible cause of the difference. Since these observations, which prove that this conjecture is untenable, were made known through the Zeitschrift für Astronomie, Prof. \textsc{Bohnenberger} has proposed a new hypothesis to explain the puzzling phenomenon. This hypothesis prompted a modification to the \textsc{Reichenbach} circle with which the above observations were made, and the effect of that modification in the observations of the most recent winter solstice seems to support the correctness of \textsc{Bohnenberger}'s hypothesis. Whether this hypothesis suffices to explain the phenomenon \emph{entirely}, however, or whether one must instead assume that \emph{several} causes act together, will be investigated elsewhere. --- Observation of the solar eclipse of 19 November 1816 at Glatz by General \textsc{von Lindener}. --- Notes on the circumstances of the late \textsc{Triesnecker}'s life by Prof. \textsc{Bürg}. --- Determinations for the Pole Star by Baron \textsc{von Lindenau}. The very valuable analysis of eight hundred observations of the Pole Star yielded several interesting results. It gave new determinations of the \SourceLeaf{613}aberration constant as $20''44861$ (probable error according to the \textsc{Gauss-Laplace} probability theory $\pm0''032$), of the nutation constant as $8''97707$ (probable error $\pm0''04421$), and of the parallax of the Pole Star as $0''14444$ (probable error $\pm0''05568$). Using the nutation constant together with the precession constant, Mr. \textsc{von Lindenau} also tried to derive the Earth's flattening and the Moon's mass by means of some formulas from the Mécanique Céleste and an equation of condition communicated to him by \textsc{Gauss}. For the former he found a value close to $\tfrac{1}{316}$. Several errors, however, crept into the formulas used and into the calculation, and Mr. \textsc{von Lindenau} himself can be expected to correct them in a suitable place. It is mere chance that the above result comes so close to the truth. If all the errors are properly corrected, the flattening comes out at about $\tfrac{1}{203}$. It must be noted, however, that the equation of condition given by \textsc{Gauss} (from which the coefficient $\tfrac{3}{5}$ must be omitted in the form printed here) is hypothetical insofar as it assumes that the layers of equal density in the Earth's spheroid are similar. --- Oppositions of Jupiter and Uranus observed in 1817 at the Seeberg Observatory by Assistant \textsc{Enke}. The latter opposition is especially important because Uranus has never before been observed so close to its node. --- On refining an already approximately known cometary orbit, by Dr. \textsc{Olbers}. An inappropriate judgment that an astronomer had made about the method proposed by \textsc{Olbers} for the stated purpose is corrected here. It is also shown how this method can be adapted to determine the elliptical orbit by adding a fourth hypothesis. This useful procedure deserves to be recommended alongside others. The author also rightly observes that it would be an unjustifiable waste of effort and time to undertake to compute the orbit as an ellipse for \emph{all} comets. This judgment, however, should deter no one with the ability to do so from determining with all due care the orbit of a comet whose observed motion clearly departs from a parabola. The highly distinguished author will surely have no objection, either, if aspiring astronomers, to hone their abilities for a task that is, after all, not within everyone's reach, perhaps even try their hand at a comet for which the result ultimately shows only a departure from the parabola confined within rather uncertain limits. --- New calculation of the secular variations in the elements of the Earth's orbit, by Prof. \textsc{Nicolai}. This very valuable work was prompted by the conjecture that the difference between the determination of the mass of Venus that Mr. \textsc{von Lindenau} obtained from the theory of Mercury and that inferred from the decrease in the obliquity of the ecliptic according to \textsc{Laplace}'s theory might arise because Laplace gives only an approximate determination of the secular variations. The author therefore recalculated the secular variations in all the elements of the Earth's orbit with the greatest care, using \textsc{Gauss}'s new method (not yet made public), which is completely rigorous. The outcome, however, did not confirm that conjecture, since the new results did not turn out to differ substantially from the older ones. --- Observation of the oppositions of Vesta, Pallas, Jupiter and Uranus in 1817 at the Mannheim Observatory by the same astronomer. --- Geometrical course of Vesta in 1817---1818 calculated by Prof. \textsc{Gerling} in Marburg, and that of Pallas by Dr. \textsc{Tittel} in Göttingen. --- An extract from a letter by Dr. \textsc{Struve} in Dorpat reports, among other things, on a trigonometric survey of Livonia that is to be undertaken. --- Price list of astronomical instruments at the institute of Mr. \textsc{von Utzschneider} in Benediktbeuern. --- The miscellaneous astronomical observations, remarks and reports that close this year's volume also contain several astronomical and geographical determinations that we cannot list individually here for lack of space.
\par\clearpage\endgroup
